\chapter{Conclusion and Future Scope}
\label{ch:conclusion}

\begin{center}
\textit{This chapter concludes the project and lists the identified future
scope together with the extensions that the architecture is structured to
support.}
\end{center}

\vspace*{\fill}
\clearpage
\section{Conclusion}

\noindent This project set out to build a single, tenant-isolated platform that
walks a college through its content-preparation and examination lifecycle, and
it delivered one: institutes are provisioned with users, roles, and
permissions; academic structure is modelled and turned into enforceable
academic scope; uploaded materials are OCR-processed by distributed workers
into reusable text; syllabus documents are analysed into confirmed
curricula; study content and question banks are generated with provenance
through an internal AI gateway; paper patterns, question papers, and
scheduled assessments are authored, published, attempted, and evaluated
server-side for objective answer formats; aggregate analytics follow; an
ungraded practice mode supports revision; and content, papers, patterns, and
results export to PDF, DOCX, and XLSX.

The engineering outcome is equally deliberate. The system is a modular
monolith with asynchronous workers, a single PostgreSQL database, strict
tenant scoping at the enforcement layers, and a security design ---
session-bound tokens, strict refresh rotation, double-submit CSRF, DB-fresh
permission evaluation from an explicit 83-key catalogue, academic scope, and a
documentation-verified audit --- that resists the access-control leaks that
plague role-name-based systems. The implementation is validated by 226
passing unit tests, eleven database-integrated suites, strict type-checking
and linting across both languages, and a recorded user-validation pass, with
no unresolved security findings.

\section{Future Scope}

\noindent The following extensions are identified and left for future work:

\begin{enumerate}
  \item \textbf{Checking-system features} --- online answer sheets (FORM),
  OMR scanning, and on-screen marking, including automated grading of
  subjective text answers, which the current evaluation contract deliberately
  records but does not grade.
  \item \textbf{Attempt mechanics} --- true per-attempt N-of-M numbering with
  maximum-attempt enforcement and auto-replenishment after selection.
  \item \textbf{Offering-level academic scope} --- introducing the designed
  \texttt{offeringId} binding so resources can additionally be restricted to a
  single class--subject offering.
  \item \textbf{Selection and deduplication quality} --- randomized selection
  within pattern buckets, similarity-based de-duplication across the bank, and
  a read path for question origin patterns.
  \item \textbf{Platform administration surface} --- a working
  \texttt{SUPER\_ADMIN} console for institute lifecycle and the OCR worker
  registry.
  \item \textbf{Operational hardening} --- distributed Redis-backed rate
  limiting, S3-backed storage behind the existing storage abstraction,
  quantitative OCR/AI quality metrics, worker observability, and web-level
  end-to-end test automation.
  \item \textbf{Academic export redesign} --- per-topic and per-resource
  exports with richer document formats.
\end{enumerate}

\noindent The architecture is structured so that each extension lands as a new
module or a role-gated surface without disturbing the tenant-isolation and
authorization guarantees that this report describes.